\begin{table}[tp]
\centering
\caption{Stochastic monotone root finding under generalized Lipschitz continuity: \href{https://github.com/k3nfalt/longtails/blob/main/LongTails/Setups/StochasticGeneralizedMonotone.lean}{StochasticGeneralizedMonotone}.}\label{tab:StochasticGeneralizedMonotone}
\parbox{\linewidth}{\footnotesize
On a finite-dimensional real Hilbert space, $F$ is continuous and monotone, with $F(x_\star)=0$, $\|x_0-x_\star\|\le R$, and
$\|F(x)-F(y)\|\le (L_0+L_1\max_{z\in[x,y]}\|F(z)\|)\|x-y\|$.
The jointly measurable oracle $g$ uses IID samples of law $p$ under the probability law $P$, with integrable values and squared errors,
$\int g(x,\omega)\,dp(\omega)=F(x)$ and $\int\|g(x,\omega)-F(x)\|^2\,dp(\omega)\le\sigma^2$.
Here $L_0,L_1,R,\sigma\ge0$, $\varepsilon>0$, and $0<\delta<1$.
The target is a measurable output $X$ satisfying $P(\|F(X)\|>\varepsilon)\le\delta$.}\par\smallskip
\begin{threeparttable}
\footnotesize
\setlength{\tabcolsep}{4pt}\renewcommand{\arraystretch}{1.25}
\begin{tabularx}{\linewidth}{@{}>{\raggedright\arraybackslash}p{0.20\linewidth}>{\raggedright\arraybackslash}p{0.32\linewidth}>{\raggedright\arraybackslash}p{0.25\linewidth}>{\raggedright\arraybackslash}X@{}}
\toprule
\textbf{Result} & \textbf{Oracle complexity} & \textbf{Novelty} & \textbf{Limitations} \\
\midrule
\textbf{(1)} \href{https://github.com/k3nfalt/longtails/blob/main/LongTails/Theorems/StochasticGeneralizedMonotone.lean}{Stochastic\allowbreak Generalized\allowbreak Monotone} \textbf{(Ours)} $(\star\star\star)$ \href{https://github.com/k3nfalt/longtails/blob/main/scripts/reviews/StochasticGeneralizedMonotone.md}{(review)}
& $\mathcal{O}\!\left(\left[1+\frac{L_0R}{\varepsilon}+(L_1R)^2+\frac{\sigma^2}{\varepsilon^2}\right]\right.$\newline
$\left.{}\cdot\ell^3[\ell+\log(1/\delta)]\right)$ stochastic calls, including batches and confidence trials.
& & \\
\midrule[1.2pt]
\textbf{(2)} Accumulative regularization \citep{sgmJiLanZhu2026v1}, Thm.~2.2; special case $L_1=0$
& $\mathcal{O}\!\left(\frac{L_0R}{\varepsilon}h+\frac{\sigma^2}{\varepsilon^2}\ell_0^3h^2\right)$ stochastic calls for $\sqrt{\mathbb{E}\|F(X)\|^2}\le\varepsilon$; $0<\varepsilon\le L_0R$.
& \textbf{(1)}: Scope extends to $L_1>0$ with polynomial dependence. At $L_1=0$, matches polynomial accuracy dependence; high-probability residual instead of expected residual, with logarithmic confidence cost.
& \textbf{(1)}: Extra logarithmic overhead at $L_1=0$; unconstrained roots only, whereas \textbf{(2)} also covers constrained variational inequalities. \\
\midrule
\textbf{(3)} Clipped Korpelevich \citep{sgmVankov2025v2}, Thm.~4.4, $p=2$
& $\mathcal{O}\!\left(1+\frac{dR}{a\sqrt{\tau}}+\frac{\sigma^2+K_1\sigma^{2\alpha}}{a^2\tau}\right)$ stochastic calls for $\mathbb{E}\operatorname{dist}^2(X,X^\star)\le\tau$; two calls per iteration.
& \textbf{(1)}: No quasi-sharpness; permits $\alpha=1$. \textbf{(3)} requires $2$-quasi-sharpness with $\mu>0$ and $0<\alpha\le\tfrac12$. Distance and residual rates are not comparable.
& \textbf{(1)}: Restricted to monotone unconstrained roots; \textbf{(3)} also covers constrained and structured non-monotone problems. \\
\midrule
\textbf{(4)} Anchored extragradient \citep{cgmRepositoryAnchored}; exact-oracle special case $\sigma=0$, $L_0>0$
& $\mathcal{O}\!\left(1+\frac{L_0R}{\varepsilon}+(1+L_1R)\right.$\newline
$\left.{}\cdot\log\!\left(1+\frac{(L_0+L_1G)R}{\varepsilon}\right)\right)$ exact calls for $\|F(X)\|\le\varepsilon$.
& \textbf{(1)}: Handles finite-variance stochastic oracles; budget needs no initial-residual bound $G$. \textbf{(4)} supplies no guarantee for $\sigma>0$.
& \textbf{(1)}: At $\sigma=0$, quadratic rather than linear $L_1R$ term, with extra logarithms; confidence-selected output instead of an anytime last iterate. \\
\bottomrule
\end{tabularx}
\begin{tablenotes}[flushleft]\footnotesize
\item[] \textbf{Abbr:} $\ell=1+\log(1+L_0R/\varepsilon+L_1R)$, $\ell_0=1+\log(1+L_0R/\varepsilon)$, $h=1+\log\ell_0$; $G\ge\|F(x_0)\|$.
For \textbf{(3)}, $\tau>0$ is the squared-distance tolerance, $X^\star$ the solution set, $R\ge\operatorname{dist}(x_0,X^\star)$, and $2$-quasi-sharpness means $\langle F(x),x-x_\star\rangle\ge\mu\operatorname{dist}^2(x,X^\star)$ for every solution $x_\star$.
Its $\alpha$-symmetric condition replaces $\|F(z)\|$ by $\|F(z)\|^\alpha$ in the segment bound above.
$C_F\ge\sup_k\mathbb{E}\|F(h_k)\|$, $a=\mu\min\{1,[2(C_F+\sigma)]^{-1}\}$, $d=\max\{4\mu,2\sqrt3(K_0+K_1+K_2)\}$,
$K_0=L_0(2^{\alpha^2/(1-\alpha)}+1)$, $K_1=L_1 2^{\alpha^2/(1-\alpha)}3^\alpha$, and $K_2=L_1^{1/(1-\alpha)}2^{\alpha^2/(1-\alpha)}3^\alpha(1-\alpha)^{\alpha/(1-\alpha)}$.
\end{tablenotes}
\end{threeparttable}
\end{table}
\clearpage
